\documentclass[main.tex]{subfiles}

\begin{document}

\section{Heater Box PLC I/O Interface}

The following tables define the PLC I/O mapping for the Heater Control Box. All signals are routed through the DB37 passthrough interface board.

\subsection{Digital Inputs (5069-IB16)}

\begin{table}[H]
\centering
\caption{Heater Box Digital Inputs}
\label{tab:heater-din}
\begin{tabularx}{\textwidth}{l l l X}
\hline
\textbf{IB16 Channel} & \textbf{DB37 Pin} & \textbf{Signal} & \textbf{Description} \\
\hline
Din0 & 3 & Pushbutton 1 & Operator input (e.g., Start / Run) \\
Din1 & 4 & Pushbutton 2 & Operator input (e.g., Stop / Acknowledge) \\
\hline
\end{tabularx}
\end{table}

\subsection{Digital Outputs (5069-OB16)}

\begin{table}[H]
\centering
\caption{Heater Box Digital Outputs}
\label{tab:heater-dout}
\begin{tabularx}{\textwidth}{l l l X}
\hline
\textbf{OB16 Channel} & \textbf{DB37 Pin} & \textbf{Signal} & \textbf{Description} \\
\hline
Dout0 & 13 & Green LED & Status indicator (e.g., Run / Heater Active) \\
Dout1 & 14 & Yellow LED & Status indicator (e.g., Warning) \\
Dout2 & 15 & Red LED & Status indicator (e.g., Fault / Alarm) \\
Dout3 & 16 & Heater Relay & Energizes the 150\,\si{\ohm} heater element \\
Dout4 & 17 & Fan & Thermal disturbance / cooling \\
\hline
\end{tabularx}
\end{table}

\subsection{Analog Inputs (5069-IY4)}

\begin{table}[H]
\centering
\caption{Heater Box Analog Inputs}
\label{tab:heater-ain}
\begin{tabularx}{\textwidth}{l l l X}
\hline
\textbf{IY4 Channel} & \textbf{DB37 Pins} & \textbf{Signal} & \textbf{Description} \\
\hline
Ain0 & 21--22 & Potentiometer & 0--10\,V setpoint from front panel pot \\
Ain1 & 23--24 & Temperature & 4--20\,mA from ProSense XTD2 transmitter \\

\hline
\end{tabularx}
\end{table}

\subsection{Analog Outputs (5069-OF4)}

\begin{table}[H]
\centering
\caption{Heater Box Analog Outputs}
\label{tab:heater-aout}
\begin{tabularx}{\textwidth}{l l l X}
\hline
\textbf{OF4 Channel} & \textbf{DB37 Pins} & \textbf{Signal} & \textbf{Description} \\
\hline
Aout0 & 29--30 & Voltmeter & 0--10\,V output driving the front panel analog voltmeter \\
\hline
\end{tabularx}
\end{table}

\end{document}
